\documentclass[../main.tex]{subfiles}

\begin{document}
\begin{landscape}

\chapter*{PHỤ LỤC}
\addcontentsline{toc}{chapter}{PHỤ LỤC}

\noindent\textbf{Bảng tổng hợp kết quả khảo sát}

\tiny
{
\renewcommand{\arraystretch}{1.5}
\begin{tabularx}{\textwidth}{|p{1cm}|p{1.50cm}|p{1.35cm}|p{1.3cm}|p{1.50cm}|p{1.55cm}|p{1.55cm}|p{1.55cm}|p{1.55cm}|p{2.50cm}|p{1.55cm}|p{1.55cm}|p{1.55cm}|}
\hline
\textbf{Dấu thời gian} & \textbf{Bạn hiện đang là?} & \textbf{Bạn đã học lập trình trong bao lâu?} & \textbf{Bạn đang học ngôn ngữ lập trình nào?} & \textbf{Bạn thường học lập trình thông qua nguồn nào?} & \textbf{Bạn thường làm gì sau khi học lý thuyết lập trình?} & \textbf{Khó khăn lớn nhất của bạn khi học lập trình là gì?} & \textbf{Khi gặp lỗi code, bạn thường làm gì?} & \textbf{Bạn có sử dụng AI để hỗ trợ học lập trình không?} & \textbf{Bạn mong muốn hệ thống AI tutor (Gia sư AI) hỗ trợ những chức năng nào?} & \textbf{Bạn có sẵn sàng sử dụng một hệ thống AI tutor để hỗ trợ học lập trình không?} & \textbf{Theo bạn, một hệ thống AI tutor (Gia sư AI) nên có chức năng gì để hỗ trợ học lập trình hiệu quả?} & \textbf{Bạn có muốn tham gia test hệ thống AI tutor trong tương lai không?} \\
\hline
3/14/2026 23:26:51 & Sinh viên ngành khác & Trên 2 năm & Java & Udemy / Coursera & Làm bài tập thực hành & Không hiểu khái niệm lập trình & Hỏi bạn bè / giảng viên & Thường xuyên & Gợi ý cách sửa lỗi & Có & Không biết & Có \\
\hline
3/14/2026 23:31:20 & Sinh viên ngành khác & Dưới 6 tháng & C++ & Youtube & Không làm gì thêm & Không biết bắt đầu từ đâu & Sử dụng AI (ChatGPT, Gemini, ...) & Thỉnh thoảng & Phân tích lỗi code, Trả lời câu hỏi khi học & Có &  & Có \\
\hline
3/14/2026 23:31:48 & Sinh viên ngành khác & 6 tháng - 1 năm & JavaScript & Tài liệu chính thức (Documentation) & Viết project nhỏ & Không biết bắt đầu từ đâu & Hỏi bạn bè / giảng viên & Thường xuyên & Giải thích khái niệm lập trình, Phân tích lỗi code, Gợi ý cách sửa lỗi, Gợi ý bài học tiếp theo, Trả lời câu hỏi khi học & Có &  & Có \\
\hline
3/14/2026 23:34:12 & Sinh viên ngành khác & Dưới 6 tháng & C++ & Youtube & Làm bài tập thực hành & Không hiểu khái niệm lập trình & Tìm lỗi trên Google & Thường xuyên & Giải thích khái niệm lập trình & Có &  & Có \\
\hline
3/15/2026 6:28:44 & Sinh viên CNTT & Trên 2 năm & Java & Youtube & Làm bài tập thực hành & Không biết bắt đầu từ đâu & Sử dụng AI (ChatGPT, Gemini, ...) & Thường xuyên & Giải thích khái niệm lập trình, Phân tích lỗi code, Gợi ý cách sửa lỗi & Có & Giúp giải thích bài toán & Có \\
\hline
3/15/2026 9:34:20 & Sinh viên CNTT & Trên 2 năm & Python & Youtube & Làm bài tập thực hành & Không có người hướng dẫn & Sử dụng AI (ChatGPT, Gemini, ...) & Thường xuyên & Giải thích khái niệm lập trình, Phân tích lỗi code, Gợi ý cách sửa lỗi, Gợi ý bài học tiếp theo, Trả lời câu hỏi khi học & Có &  & Có \\
\hline
3/15/2026 10:02:25 & Sinh viên ngành khác & 6 tháng - 1 năm & JavaScript & Hỏi đáp AI (ChatGPT, Gemini, ...) & Làm bài tập thực hành & Không biết sửa lỗi code & Sử dụng AI (ChatGPT, Gemini, ...) & Thường xuyên & Phân tích lỗi code, Gợi ý cách sửa lỗi & Có thể &  & Có \\
\hline
3/15/2026 13:03:31 & Sinh viên CNTT & Trên 2 năm & Java & Youtube & Làm bài tập thực hành & Không có người hướng dẫn & Sử dụng AI (ChatGPT, Gemini, ...) & Thường xuyên & Phân tích lỗi code & Có thể &  & Có \\
\hline
3/15/2026 13:43:22 & Người tự học lập trình & Dưới 6 tháng & C++ & Hỏi đáp AI (ChatGPT, Gemini, ...) & Không làm gì thêm & Không có người hướng dẫn & Sử dụng AI (ChatGPT, Gemini, ...) & Chưa từng & Giải thích khái niệm lập trình, Phân tích lỗi code, Gợi ý cách sửa lỗi, Gợi ý bài học tiếp theo, Trả lời câu hỏi khi học & Không & Gửi link & Có \\
\hline
3/15/2026 14:52:43 & Sinh viên CNTT & Dưới 6 tháng & JavaScript & Hỏi đáp AI (ChatGPT, Gemini, ...) & Làm bài tập thực hành & Không có người hướng dẫn & Sử dụng AI (ChatGPT, Gemini, ...) & Thường xuyên & Phân tích lỗi code, Gợi ý bài học tiếp theo, Trả lời câu hỏi khi học & Có thể & giải thích code chi tiết, phát hiện và sửa lỗi lập trình, đưa ra bài tập thực hành phù hợp với trình độ, gợi ý cách giải khi người học gặp khó khăn, và theo dõi tiến độ học tập & Có \\
\hline
3/15/2026 14:55:11 & Sinh viên CNTT & 6 tháng - 1 năm & C++ & Khóa học ở trường & Làm bài tập thực hành & Không biết bắt đầu từ đâu & Tìm lỗi trên Google & Thường xuyên & Phân tích lỗi code & Có &  & Có \\
\hline
3/15/2026 14:55:16 & Sinh viên CNTT & Dưới 6 tháng & Python & Khóa học ở trường & Làm bài tập thực hành & Không biết bắt đầu từ đâu & Sử dụng AI (ChatGPT, Gemini, ...) & Thường xuyên & Giải thích khái niệm lập trình, Phân tích lỗi code, Gợi ý cách sửa lỗi, Gợi ý bài học tiếp theo, Trả lời câu hỏi khi học & Có &  & Có \\
\hline
3/15/2026 14:58:55 & Sinh viên CNTT & Trên 2 năm & Python & Tài liệu chính thức (Documentation) & Làm bài tập thực hành & Thiếu bài tập thực hành & Sử dụng AI (ChatGPT, Gemini, ...) & Thường xuyên & Giải thích khái niệm lập trình, Phân tích lỗi code, Gợi ý cách sửa lỗi, Gợi ý bài học tiếp theo & Có thể &  & Có \\
\hline
3/15/2026 14:59:07 & Sinh viên CNTT & Trên 2 năm & C++ & Youtube & Làm bài tập thực hành & Không hiểu khái niệm lập trình & Tìm lỗi trên Google & Thường xuyên & Giải thích khái niệm lập trình, Phân tích lỗi code, Gợi ý cách sửa lỗi & Có &  & Có \\
\hline
3/15/2026 15:12:11 & Sinh viên CNTT & Dưới 6 tháng & Python & Hỏi đáp AI (ChatGPT, Gemini, ...) & Không làm gì thêm & Không có người hướng dẫn & Sử dụng AI (ChatGPT, Gemini, ...) & Thường xuyên & Phân tích lỗi code, Gợi ý cách sửa lỗi, Trả lời câu hỏi khi học & Có &  & Có \\
\hline
3/15/2026 15:13:11 & Sinh viên CNTT & 1 - 2 năm & C++ & Khóa học ở trường & Làm bài tập thực hành & Không biết bắt đầu từ đâu & Sử dụng AI (ChatGPT, Gemini, ...) & Thường xuyên & Giải thích khái niệm lập trình, Phân tích lỗi code, Gợi ý cách sửa lỗi, Gợi ý bài học tiếp theo, Trả lời câu hỏi khi học & Có &  & Có \\
\hline
3/15/2026 15:18:48 & Sinh viên CNTT & Trên 2 năm & dart & Hỏi đáp AI (ChatGPT, Gemini, ...) & Đọc thêm tài liệu & Thiếu bài tập thực hành & Sử dụng AI (ChatGPT, Gemini, ...) & Thường xuyên & Trả lời câu hỏi khi học & Có thể & nhắc nhở học tập và đưa cho người học lộ trình học tập đầy đủ & Có \\
\hline
3/15/2026 15:20:23 & Sinh viên CNTT & Trên 2 năm & Python & Khóa học ở trường & Không làm gì thêm & Không có người hướng dẫn & Sử dụng AI (ChatGPT, Gemini, ...) & Thường xuyên & Gợi ý cách sửa lỗi & Không & Đố biết & Có \\
\hline
3/15/2026 15:39:49 & Sinh viên CNTT & Trên 2 năm & Java & Khóa học ở trường & Viết project nhỏ & Thiếu bài tập thực hành & Sử dụng AI (ChatGPT, Gemini, ...) & Thỉnh thoảng & Phân tích lỗi code, Gợi ý cách sửa lỗi, Gợi ý bài học tiếp theo & Có thể &  & Có \\
\hline
3/15/2026 15:53:18 & Sinh viên CNTT & Trên 2 năm & Java & Hỏi đáp AI (ChatGPT, Gemini, ...) & Làm bài tập thực hành & Không hiểu khái niệm lập trình & Sử dụng AI (ChatGPT, Gemini, ...) & Thường xuyên & Giải thích khái niệm lập trình & Có &  & Có \\
\hline
3/15/2026 16:07:45 & Sinh viên CNTT & Trên 2 năm & Python & Youtube & Làm bài tập thực hành & Không biết bắt đầu từ đâu & Sử dụng AI (ChatGPT, Gemini, ...) & Thường xuyên & Giải thích khái niệm lập trình & Có &  & Có \\
\hline
3/15/2026 16:22:46 & Sinh viên CNTT & Trên 2 năm & Python & Tài liệu chính thức (Documentation) & Làm bài tập thực hành & Không biết bắt đầu từ đâu & Sử dụng AI (ChatGPT, Gemini, ...) & Thường xuyên & Giải thích khái niệm lập trình, Phân tích lỗi code, Gợi ý cách sửa lỗi, Trả lời câu hỏi khi học & Có thể &  & Có \\
\hline
3/15/2026 18:11:49 & Sinh viên CNTT & Trên 2 năm & JavaScript & Hỏi đáp AI (ChatGPT, Gemini, ...) & Làm bài tập thực hành & Không hiểu khái niệm lập trình & Sử dụng AI (ChatGPT, Gemini, ...) & Thỉnh thoảng & Phân tích lỗi code & Có &  & Có \\
\hline
3/15/2026 19:38:46 & Sinh viên CNTT & Trên 2 năm & JavaScript & Youtube & Viết project nhỏ & Không biết bắt đầu từ đâu & Sử dụng AI (ChatGPT, Gemini, ...) & Thường xuyên & Phân tích lỗi code, Gợi ý cách sửa lỗi, Gợi ý bài học tiếp theo & Có thể & Tư vấn lộ trình, lên kế hoạch chi tiết và bài học cho từng giai đoạn để người sử dụng nắm bắt được quá trình học tập và mục tiêu cụ thể & Có \\
\hline
3/15/2026 19:50:36 & Sinh viên CNTT & Trên 2 năm & C++ & Tài liệu chính thức (Documentation) & Viết project nhỏ & Không hiểu khái niệm lập trình & Tìm lỗi trên Google & Thường xuyên & Giải thích khái niệm lập trình & Có & Tuỳ & Có \\
\hline
3/15/2026 23:50:13 & Sinh viên CNTT & 6 tháng - 1 năm & C++ & Khóa học ở trường & Đọc thêm tài liệu & Không biết sửa lỗi code & Hỏi bạn bè / giảng viên & Thỉnh thoảng & Giải thích khái niệm lập trình & Có &  & Có \\
\hline
3/15/2026 23:57:25 & Sinh viên CNTT & 6 tháng - 1 năm & C++ & Khóa học ở trường & Làm bài tập thực hành & Không biết sửa lỗi code & Hỏi bạn bè / giảng viên & Thường xuyên & Phân tích lỗi code, Gợi ý bài học tiếp theo & Có thể & Sửa lỗi và đưa ra lịch trình & Có \\
\hline
3/16/2026 0:00:11 & Sinh viên CNTT & Dưới 6 tháng & C++ & Khóa học ở trường & Làm bài tập thực hành & Không biết bắt đầu từ đâu & Hỏi bạn bè / giảng viên & Thường xuyên & Phân tích lỗi code, Gợi ý cách sửa lỗi, Gợi ý bài học tiếp theo & Có & giải thích kĩ về cách giải 1 bài tập nào đó & Có \\
\hline
3/16/2026 0:42:33 & Sinh viên CNTT & Dưới 6 tháng & C & Youtube & Đọc thêm tài liệu & Không biết bắt đầu từ đâu & Sử dụng AI (ChatGPT, Gemini, ...) & Thường xuyên & Giải thích khái niệm lập trình, Phân tích lỗi code, Gợi ý cách sửa lỗi, Gợi ý bài học tiếp theo, Trả lời câu hỏi khi học & Có thể & Tôi thấy nếu như là 1 giá sư thì phải có chức năng giúp sinh viên cải thiện bài tập và thêm bài tập cho sinh viên để sinh viên hiểu hơn về lập trình . Giải thích cẵn kế chỉ tiết cho Sinh viên & Có \\
\hline
3/16/2026 2:47:43 & Sinh viên CNTT & Dưới 6 tháng & C++ & Khóa học ở trường & Làm bài tập thực hành & Không biết bắt đầu từ đâu & Sử dụng AI (ChatGPT, Gemini, ...) & Thường xuyên & Giải thích khái niệm lập trình & Có &  & Có \\
\hline
3/16/2026 5:22:20 & Sinh viên CNTT & Dưới 6 tháng & C++ & Khóa học ở trường & Làm bài tập thực hành & Không biết sửa lỗi code & Sử dụng AI (ChatGPT, Gemini, ...) & Thường xuyên & Phân tích lỗi code & Có & Nhiều tnang & Có \\
\hline
3/16/2026 5:37:17 & Sinh viên CNTT & Dưới 6 tháng & C++ & Khóa học ở trường & Viết project nhỏ & Không biết bắt đầu từ đâu & Tìm lỗi trên Google & Thỉnh thoảng & Giải thích khái niệm lập trình, Phân tích lỗi code, Gợi ý cách sửa lỗi & Có thể &  & Không \\
\hline
3/16/2026 7:42:47 & Sinh viên CNTT & Dưới 6 tháng & C++ & Khóa học ở trường & Làm bài tập thực hành & Không biết bắt đầu từ đâu & Sử dụng AI (ChatGPT, Gemini, ...) & Thỉnh thoảng & Giải thích khái niệm lập trình, Phân tích lỗi code, Gợi ý cách sửa lỗi, Gợi ý bài học tiếp theo & Có thể &  & Không \\
\hline
3/16/2026 7:43:14 & Sinh viên CNTT & 6 tháng - 1 năm & C++ & Hỏi đáp AI (ChatGPT, Gemini, ...) & Làm bài tập thực hành & Không biết sửa lỗi code & Sử dụng AI (ChatGPT, Gemini, ...) & Thỉnh thoảng & Phân tích lỗi code, Gợi ý cách sửa lỗi & Có &  & Có \\
\hline
3/16/2026 9:04:58 & Sinh viên CNTT & 6 tháng - 1 năm & C++ & Hỏi đáp AI (ChatGPT, Gemini, ...) & Làm bài tập thực hành & Không biết sửa lỗi code & Hỏi bạn bè / giảng viên & Thỉnh thoảng & Phân tích lỗi code & Có &  & Không \\
\hline
3/16/2026 10:59:53 & Sinh viên CNTT & Dưới 6 tháng & C++ & Khóa học ở trường & Làm bài tập thực hành & Không biết bắt đầu từ đâu & Sử dụng AI (ChatGPT, Gemini, ...) & Thỉnh thoảng & Giải thích khái niệm lập trình, Phân tích lỗi code, Gợi ý cách sửa lỗi, Gợi ý bài học tiếp theo, Trả lời câu hỏi khi học & Có &  & Không \\
\hline
3/16/2026 11:14:38 & Sinh viên CNTT & Dưới 6 tháng & C++ & Khóa học ở trường & Đọc thêm tài liệu & Không biết sửa lỗi code & Hỏi bạn bè / giảng viên & Thỉnh thoảng & Phân tích lỗi code & Có thể &  & Có \\
\hline
3/17/2026 8:43:17 & Sinh viên CNTT & Dưới 6 tháng & C++ & Youtube & Làm bài tập thực hành & Không biết sửa lỗi code & Sử dụng AI (ChatGPT, Gemini, ...) & Thường xuyên & Phân tích lỗi code & Có thể & có & Có \\
\hline
3/17/2026 9:29:15 & Sinh viên CNTT & Dưới 6 tháng & C++ & Khóa học ở trường & Làm bài tập thực hành & Không biết bắt đầu từ đâu & Tìm lỗi trên Google & Thỉnh thoảng & Giải thích khái niệm lập trình, Phân tích lỗi code, Gợi ý cách sửa lỗi, Gợi ý bài học tiếp theo & Có thể &  & Có \\
\hline
3/17/2026 11:05:16 & Sinh viên CNTT & Dưới 6 tháng & C++ & Youtube & Làm bài tập thực hành & Không biết bắt đầu từ đâu & Sử dụng AI (ChatGPT, Gemini, ...) & Thường xuyên & Giải thích khái niệm lập trình, Phân tích lỗi code, Gợi ý cách sửa lỗi, Gợi ý bài học tiếp theo & Có &  & Có \\
\hline
3/17/2026 11:18:16 & Sinh viên CNTT & Dưới 6 tháng & C++ & Tài liệu chính thức (Documentation) & Đọc thêm tài liệu & Không biết sửa lỗi code & Sử dụng AI (ChatGPT, Gemini, ...) & Thường xuyên & Giải thích khái niệm lập trình, Phân tích lỗi code, Gợi ý cách sửa lỗi, Gợi ý bài học tiếp theo, Trả lời câu hỏi khi học & Có & Giải đáp nhanh đúng và chính xác & Có \\
\hline
3/17/2026 12:56:30 & Sinh viên CNTT & Dưới 6 tháng & C++ & Khóa học ở trường & Làm bài tập thực hành & Không biết sửa lỗi code & Sử dụng AI (ChatGPT, Gemini, ...) & Thường xuyên & Phân tích lỗi code & Có thể &  & Không \\
\hline
3/17/2026 15:50:16 & Sinh viên CNTT & Dưới 6 tháng & C++ & Khóa học ở trường & Làm bài tập thực hành & Không biết bắt đầu từ đâu & Sử dụng AI (ChatGPT, Gemini, ...) & Hiếm khi & Giải thích khái niệm lập trình & Có &  & Có \\
\hline
3/17/2026 16:15:10 & Sinh viên CNTT & Dưới 6 tháng & C++ & Khóa học ở trường & Làm bài tập thực hành & Không biết sửa lỗi code & Sử dụng AI (ChatGPT, Gemini, ...) & Thường xuyên & Giải thích khái niệm lập trình, Phân tích lỗi code & Có thể &  & Có \\
\hline
3/17/2026 16:57:56 & Sinh viên CNTT & Dưới 6 tháng & C++ & Youtube & Làm bài tập thực hành & Không có người hướng dẫn & Sử dụng AI (ChatGPT, Gemini, ...) & Thỉnh thoảng & Giải thích khái niệm lập trình & Có & Kĩ năng giải thích & Có \\
\hline
3/18/2026 10:08:21 & Sinh viên CNTT & Dưới 6 tháng & C++ & Hỏi đáp AI (ChatGPT, Gemini, ...) & Làm bài tập thực hành & Không biết sửa lỗi code & Tìm trên StackOverflow & Thường xuyên & Gợi ý cách sửa lỗi & Có &  & Có \\
\hline
3/18/2026 10:33:33 & Sinh viên CNTT & Dưới 6 tháng & C++ & Khóa học ở trường & Làm bài tập thực hành & Không hiểu khái niệm lập trình & Hỏi bạn bè / giảng viên & Hiếm khi & Giải thích khái niệm lập trình & Không &  & Có \\
\hline
\end{tabularx}
}


\end{landscape}
\end{document}